\section{An explicit capped-geometric Gaussian submodel}
\label{sec:explicit-capped-gaussian}

We now remove both abstract constants in
Theorem~\ref{thm:hsmm-end-to-end} for one named family.  Fix $D\ge2$, known
$\sigma>0$, and
\[
 \theta=(a,p_1,p_2)\in
 \Theta_{\rm cg}:=[a_-,a_+]\times[p_-,p_+]^2,
 \quad 0<a_-<a_+,\quad0<p_-<p_+<1.
\]
Let $D_i=\min(G_i,D)$ with $G_i\sim{\rm Geom}(p_i)$, so that
\[
 f_i(d)=p_i(1-p_i)^{d-1}\quad(d<D),\qquad
 f_i(D)=(1-p_i)^{D-1}.
\]
This is capping, not the geometric law conditioned on $G_i\le D$.  Put
$Y_t=2\mathbf1_{\{S_t=2\}}-1$ and
\[
 X_t=aY_t+\sigma Z_t,\qquad Z_t\stackrel{\rm iid}{\sim}N(0,1).
\]
Choose $\omega$ such that $0<\omega a_+<\pi/2$, define
$H(x)=2\mathbf1_{\{x>0\}}-1$, and use the fixed bounded feature vector
\begin{equation}
 \phi(x_0,x_1,x_2)=
 \bigl(\cos(\omega x_0),H(x_0),H(x_0)H(x_1),
 H(x_0)H(x_2),H(x_0)H(x_1)H(x_2)\bigr).
 \label{eq:cg-feature}
\end{equation}
Let $z=(z_0,\ldots,z_4)=\E_\theta\phi(X_0,X_1,X_2)$ and
$\delta(a)=2\Phi_{\rm N}(a/\sigma)-1$.  Then the exact inverse is
\begin{align}
 a&=\omega^{-1}\arccos\{e^{\sigma^2\omega^2/2}z_0\},
 &\alpha:=\E Y_t&=z_1/\delta(a),\nonumber\\
 r_1:=\E(Y_tY_{t+1})&=z_2/\delta(a)^2,
 &r_2:=\E(Y_tY_{t+2})&=z_3/\delta(a)^2,\nonumber\\
 t_3:=\E(Y_tY_{t+1}Y_{t+2})&=z_4/\delta(a)^3,\nonumber\\
 p_1&=\frac{1+\alpha-2r_1+r_2-t_3}{2(1-r_1)},
 &p_2&=\frac{1-\alpha-2r_1+r_2+t_3}{2(1-r_1)}.
 \label{eq:cg-inverse}
\end{align}
Indeed, the directed boundary rate is
$c_\theta=(1-r_1)/4=(\tau_1+\tau_2)^{-1}$, while
$\Pp(+,-,+)=c_\theta p_1$ and
$\Pp(-,+,-)=c_\theta p_2$; expansion of the binary indicators gives
\eqref{eq:cg-inverse}.

To state a global inverse margin, define
\begin{align}
 \bar\tau_D&=\frac{1-(1-p_-)^D}{p_-},
 &\delta_-&=2\Phi_{\rm N}(a_-/\sigma)-1,\nonumber\\
 A_a&=\frac{e^{\sigma^2\omega^2/2}}
 {\omega\sin(\omega a_-)},
 &B_\delta&=\frac{\sqrt{2/\pi}}{\sigma},\nonumber\\
 L_k&=\delta_-^{-k}+kB_\delta A_a\delta_-^{-(k+1)},\quad k=1,2,3,
 &K_p&=\frac{\bar\tau_D}{4}(L_1+3L_2+L_3),\nonumber\\
 K_{\rm id,cg}&=\sqrt{A_a^2+2K_p^2},
 &\underline\kappa_{\rm id,cg}&=K_{\rm id,cg}^{-1}.
 \label{eq:cg-id-constants}
\end{align}

\begin{proposition}[Calculated identification margin]
\label{prop:cg-identification}
For all $\theta,\vartheta\in\Theta_{\rm cg}$,
\begin{equation}
 \|\theta-\vartheta\|_2
 \le K_{\rm id,cg}
 \|\E_\theta\phi-\E_\vartheta\phi\|_\infty.
 \label{eq:cg-global-inverse}
\end{equation}
Thus $\kappa_{\rm id}$ in \eqref{eq:global-inverse-margin} may be set equal
to the explicit value $\underline\kappa_{\rm id,cg}$ for this feature map.
\end{proposition}

\begin{proof}
The derivative of the first coordinate of \eqref{eq:cg-inverse} is bounded by
$A_a$.  Since $|\delta'(a)|\le B_\delta$, a ratio of an observed moment by
$\delta(a)^k$ changes by at most $L_k$ times the sup-norm error in $z$.
Moreover, $1-r_1=4c_\theta\ge2/\bar\tau_D$.  The absolute derivatives of
$p_i$ with respect to $\alpha,r_2,t_3$ are at most $\bar\tau_D/4$, and
$|\partial_{r_1}p_i|=|p_i-1|/(1-r_1)\le\bar\tau_D/2$.  Hence
$|p_i-p_i'|\le K_p\|z-z'\|_\infty$.  Combining the three coordinate bounds
in Euclidean norm proves \eqref{eq:cg-global-inverse}.
\end{proof}

The constant in \eqref{eq:cg-id-constants} is a certified global margin, not
a claim that every triangle inequality is sharp.  At a specified parameter,
the exact local margin is
\[
 \kappa_{\rm id,cg}^{\rm loc}(\theta)
 =\|D\mathcal I(\Psi(\theta))\|_{\infty\to2}^{-1}
 =\left\{
 \max_{s\in\{-1,1\}^5}\|D\mathcal I(\Psi(\theta))s\|_2
 \right\}^{-1},
\]
where $\mathcal I$ is the inverse in \eqref{eq:cg-inverse}.  Its full Jacobian
is displayed in the accompanying derivation.

For the mixing constant, augment by the current age:
$Z_t=(S_t,A_t)\in\{1,2\}\times\{1,\ldots,D\}$.  For $r<D$,
\[
 Q((i,r),(i,r+1))=1-p_i,\qquad
 Q((i,r),(3-i,1))=p_i,
\]
and $Q((i,D),(3-i,1))=1$.  Set
\begin{equation}
 a_*=\min\{p_-,1-p_+\},\qquad
 \varepsilon_*=2a_*^3,\qquad
 K_*=\left\lceil\frac{\log2}{-\log(1-\varepsilon_*)}\right\rceil.
 \label{eq:cg-mix-constants}
\end{equation}
From any augmented state, the path consisting of a switch, a continuation,
and a switch, and the path of three consecutive switches, reach $(1,1)$ and
$(2,1)$ in an order that depends on the current state.  Forced exits only
increase their probabilities.
Writing $\nu_*=\{\delta_{(1,1)}+\delta_{(2,1)}\}/2$, therefore
\begin{equation}
 Q^3(z,\cdot)\ge\varepsilon_*\nu_*(\cdot),\qquad
 \sup_{z,z'}d_{\rm TV}\{Q^t(z,\cdot),Q^t(z',\cdot)\}
 \le(1-\varepsilon_*)^{\lfloor t/3\rfloor}.
 \label{eq:cg-doeblin}
\end{equation}
For the length-three window chain, two extra shifts remove the initial
coordinates.  In the notation of Definition~1.4 of the corrected version of
Paulin~\cite{Paulin2015}, the preceding pairwise total-variation bound and the
definition of $K_*$ give
\[
 \tau_{W^{(2)}}(1/2)\le 2+3K_*.
\]
Applying the first inequality in Proposition~3.4, equation~(3.9), with
$\epsilon=1/2$ therefore yields
\[
 \gamma_{\rm ps}(W^{(2)})
 \ge \frac{1-1/2}{\tau_{W^{(2)}}(1/2)}
 \ge \frac{1}{2(2+3K_*)}.
\]
Thus the calculated uniform bound is
\begin{equation}
 \boxed{
 \underline g_3=
 \{2(2+3K_*)\}^{-1}\le\gamma_{\rm ps}(W^{(2)}).}
 \label{eq:cg-window-gap}
\end{equation}
Pointwise sharper values follow directly from the $2D\times2D$ matrix $Q$.
If $D_\eta$ is its stationary diagonal matrix, then
\begin{equation}
 \gamma_{\rm ps}(Q)=\max_{k\ge1}
 \frac{1-s_2^2(D_\eta^{1/2}Q^kD_\eta^{-1/2})}{k},
 \label{eq:cg-exact-pseudogap}
\end{equation}
where $s_2$ is the second singular value.  This is an exact finite-matrix
calculation; \eqref{eq:cg-window-gap} is the uniform analytic lower bound and
is not presented as optimal.

Finally, let $M=T-2$ and $y=x+\log5$.  Every coordinate $f$ of
\eqref{eq:cg-feature} satisfies $|f|\le1$, and hence
$|f-\E_\pi f|\le2$ and $V_f\le1$.  Since the observation process is
stationary, its length-three window chain is taken in its stationary
distribution.  Apply Theorem~3.11, equation~(3.25), of the corrected version
of Paulin~\cite{Paulin2015} to each coordinate with $n=M=T-2$,
$\gamma_{\rm ps}\ge\underline g_3$, $C\le2$, and $V_f\le1$.  That equation
controls the sum
$S=\sum_{t=1}^{M}f(W_t^{(2)})$.  Substituting the sum threshold $t=Mr$ and
requiring the exponent to be at least $y$ gives
\[
 M^2\underline g_3 r^2
 \ge y\{8(M+1/\underline g_3)+40Mr\}.
\]
Solving this quadratic inequality for $r$ and dividing the sum threshold by
$M$ gives
\begin{equation}
 r_T(x)=\frac{20y+
 \sqrt{400y^2+8y\underline g_3(M+1/\underline g_3)}}
 {M\underline g_3}.
 \label{eq:cg-explicit-radius}
\end{equation}
Applying the resulting bound to each of the five coordinates and taking a
union bound gives a factor $5$; because $y=x+\log5$, this gives
\[
 \Pp_\theta\{\|\widehat z_T-z\|_\infty>r_T(x)\}\le2e^{-x}.
\]
For the global minimum-distance estimator over $\Theta_{\rm cg}$,
Proposition~\ref{prop:cg-identification} therefore yields
\begin{equation}
 \boxed{
 \Pp_\theta\{\|\widehat\theta_T-\theta\|_2>
 2K_{\rm id,cg}r_T(x)\}\le2e^{-x}.}
 \label{eq:cg-end-to-end}
\end{equation}
The detailed proof, exact local Jacobian, and numerical matrix audit are in
\path{step1_explicit_identification_mixing.tex} and
\path{step1_explicit_constants.py}.
